\chapter{Backward reconstruction and conditional quadratic moments}\label{ch14:chapter}

This section proves three estimates on the coarse-to-fine reconstruction of the fine field under the conditioned Wilson measure. They are used in \cref{ch23:chapter,ch24:chapter} to control the heat evolution of the reconstructed field and its quadratic Picard term, and also in \cref{ch25:chapter,ch28:chapter}. Throughout, the torus, the coarsest mesh $H$, the continuum coupling $g_3>0$, and a time horizon $T_0$ are fixed, and all estimates are uniform in the number $k$ of renormalization group steps between the coarsest and the finest lattice.

\paragraph{The purpose of the three estimates.}
Recall from \cref{s3:fixed} that the fine field is reconstructed from the coarsest field one renormalization group step at a time, from coarse to fine. At the $j$th step, the logarithm $\theta_j$ of the reconstructed field on the lattice of mesh $a_j$ is obtained from the logarithm $\theta_{j-1}$ on the lattice of mesh $a_{j-1}=Ba_j$ and from the normalized fluctuation variables $w_j$ of that step. Let $A_j=a_j^{-1}\theta_j$ be the lattice potential, and let $C_j$ be the zero-fluctuation reconstruction in the same units, that is, the reconstruction with $w_j=0$. We write $\overline P$, $\overline Q$, and $\overline R$ for the linear, quadratic, and cubic parts of the zero-fluctuation reconstruction defined in \cref{ch14:lem:signed-physical-columns}, so that $C_j=a_{j-1}^{-1}[\overline P\theta_{j-1}+\overline Q(\theta_{j-1},\theta_{j-1})+\overline R(\theta_{j-1})]$. Viewing all terms as functions that are constant on lattice cells, we obtain the exact decomposition
\begin{equation}\label{ch14:eq:increment-decomposition}
 A_j-E_{a_{j-1}}A_{j-1}
 =\bigl(A_j-\E[A_j\mid\mathcal F_{j-1}]\bigr)+m_j+D_dA_{j-1}
 +a_{j-1}^{-1}\overline Q(\theta_{j-1},\theta_{j-1})+a_{j-1}^{-1}\overline R(\theta_{j-1}).
\end{equation}
Here $E_{a_{j-1}}$ is the cell embedding on the coarser lattice, and $\mathcal F_{j-1}$ is generated by the coarsest field and the fluctuation variables of the first $j-1$ steps. The term $m_j=\E[A_j\mid\mathcal F_{j-1}]-C_j$ is the conditional mean correction (in the notation of \cref{ch23:chapter,ch30:chapter}, $m_j=m_j^{\mathrm{ph}}$), and $D_dA_{j-1}$ is the difference between the linearized reconstruction of $A_{j-1}$ and its cell embedding; the precise definitions are given below. The first term is controlled by \cref{ch13:thm:centered-shell}. \Cref{ch14:lem:backward-caps-and-means} gives a deterministic bound on $\theta_j$, the reconstruction amplitude bound, which tends to zero with the mesh. It also bounds $m_j$ and shows that the sum over $j$ of $A_j-C_j$, which is the sum of the first two terms, is small in a negative-regularity norm. \Cref{ch14:lem:signed-physical-columns} shows that each column of the linear operator $D_d$ has zero integral and is localized at the scale $a_{j-1}$, so that heat smoothing at time $t$ gains the factor $a_{j-1}/\sqrt t$. \Cref{ch14:thm:quadratic-ancestor-return} controls the quadratic term. The cubic term is controlled by the reconstruction amplitude bound on $\theta_{j-1}$.

\paragraph{Inputs and attribution.}
The deterministic inputs are the reconstruction maps of \cref{ch12:thm:replacement-family}, with the representation of the conditioned Wilson measure in \cref{ch12:lem:backward-law}, the preservation of the conditional kernels in \cref{ch12:thm:representative-prefixes}, and the derivative bounds of \cref{ch12:lem:deterministic-prolongation-reserve}. The zero-fluctuation reconstruction is built from the classical background fields of Ba{\l}aban \cite{Balaban1985Variational}, and the averaging is a version of Ba{\l}aban's averaging operation \cite{Balaban1985Averaging}. The probabilistic inputs are the bounds on conditional means and covariances of \cref{ch10:thm:actual-means-and-scores}, the concentration theorem \cref{ch02:thm:reflecting-concentration}, and the results of \cref{ch13:chapter}. The concentration theorem rests on the convexity methods of Brascamp and Lieb \cite{BrascampLieb1976}, Gross \cite{Gross1975}, and Bakry and \'Emery \cite{BakryEmery1985}. The gain from a column of zero integral is the standard cancellation of a localized kernel with vanishing mass, and the maximal inequalities are elementary chaining in the sense of Dudley \cite{Dudley1967} (see also \cite{Dirksen2015}), in the form of \cref{ch13:lem:coefficient-metric-chaining}.

\section{Setting}\label{ch14:sec:setting}

\paragraph{Order of the parameters.}
We use the constructions of \cref{ch12:chapter} with the following order of choices. The constant $C_{\mathrm{CP}}\ge1$ of \cref{ch12:thm:replacement-family}, which bounds the first three derivatives in the coarse field and the first spatial difference of the zero-fluctuation reconstruction, is fixed before $J$. Before $J$ is chosen, it may be enlarged to include finitely many further bounds of the same kind, as in the proofs below. Then $J$ is chosen so that $\qanc=C_{\mathrm{CP}}^2/B<1$, where $B=M^J$; \cref{ch12:thm:replacement-family} allows $\qanc\le1/4$. The larger analytic domains and their coefficient bounds, which may depend on $J$, are fixed next. Then the radii are chosen as in \cref{ch12:lem:deterministic-prolongation-reserve},
\[
 0<r<r_{\mathrm{eff}}/4,\qquad\frac{\beta_Jr_{\mathrm{eff}}}{1-\alpha_*}<\frac18,\qquad\alpha_*=C_{\mathrm{CP}}/B,
\]
and the bounds on the coupling-scaled fluctuation inputs and on the coupling are chosen so that every reconstructed field lies in the working domain of that lemma. Finally, the upper bound $g_*$ on the couplings is fixed. It includes the bound of \cref{ch13:thm:centered-shell} with $R_{\mathrm{supp}}=r$, it ensures that $p_j\ge1$, and it makes $t_jp_j^3$ as small as required in \cref{ch10:thm:actual-means-and-scores}, in the notation introduced next.

\paragraph{The coarse-to-fine reconstruction.}
Fix $L>0$, $T_0>0$, a coarsest mesh $H=L/n_0$ with $n_0\in4\N$, and $g_3>0$ such that $t_1\le g_*$, in the notation introduced below; here $\gamma=g_3^2L$. We also fix the Borel choice of gauge of \cref{ch12:lem:backward-law}, which carries every stopping field in $D_r$ into $B_r$ and does not depend on $k$, and the reconstruction maps $\RepMap$ of \cref{ch12:thm:replacement-family}. We write $\widehat\Phi_{d,j}$ for their logarithmic form at depth $d$, at the $j$th step. For $k\ge1$, let $a_j=HB^{-j}$ and $d_j^k=J(k-j)$ for $0\le j\le k$. For $1\le j\le k$, we define
\[
 g_j=g_3\sqrt{a_j},\qquad g_{j,s}=M^{s/2}g_j\quad(0\le s<J),\qquad t_j=M^{(J-1)/2}g_j,\qquad p_j=A_{\mathrm w}\log(1/t_j).
\]
Since $t_j=t_1B^{-(j-1)/2}$, we have $t_j\le t_1\le g_*$ at every fluctuation step. Here $g_j^2$ is the lattice coupling at mesh $a_j$, $g_{j,s}$ is the coupling at the intermediate scale $s$ of the $j$th step, and $t_j$ is the largest of these couplings. Further, $p_j$ is the radius of the fluctuation cutoff domain of the $j$th step, and $d_j^k$ is the number of averaging steps that precede it, counted from the finest lattice. The reconstruction is
\begin{equation}\label{ch14:eq:backward-chain}
 \theta_j=\widehat\Phi_{d_j^k,j}\bigl(\theta_{j-1},(g_{j,s}w_{j,s})_{s=0}^{J-1}\bigr),\qquad
 \mathcal F_j=\sigma(\theta_0,w_1,\dots,w_j)\qquad(1\le j\le k),
\end{equation}
where $\theta_0$ is the logarithm of the stopping field in the chosen gauge, and $w_j=(w_{j,s})_s$ are the normalized fluctuation variables of the $j$th step. The meshes and depths are
\[
 (\theta_0;\,H;\,Jk)\longrightarrow(\theta_1;\,H/B;\,J(k-1))\longrightarrow\cdots\longrightarrow(\theta_k;\,H/B^k;\,0).
\]

\paragraph{The measure.}
Let $N=n_0B^k$, and let $\nu=\nu_{N,n_0,\gamma}$ be the conditioned Wilson measure. By \cref{ch12:lem:backward-law}, the gauge orbits of the fields $(e^{\theta_0},\dots,e^{\theta_k})$ have exactly the law of the averaged configurations under $\nu$. The conditional law of $w_j$ given $\mathcal F_{j-1}$ is the conditional fluctuation kernel $k(e^{\theta_{j-1}},dw)$ of \cref{ch10:thm:complete-white-law} at the coarse field $e^{\theta_{j-1}}$, and the reconstruction maps preserve it (\cref{ch12:thm:representative-prefixes}). This kernel has a density $Z^{-1}\one_\Omega e^{-V}$ on a fixed product $\Omega$ of three-dimensional balls of radius $p_j$, where $V$ is smooth on a neighborhood of the closed product and $D_w^2V\ge I/2$ there. It is not a product of independent Gaussian measures. All conditional expectations below are taken under $\nu$. The marginal law of $\theta_{j-1}$ under $\nu$ includes the restrictions at all coarser steps, and it is in general not a stopped measure. Once $\mathcal F_{j-1}$ is fixed, these restrictions are fixed factors, and they cancel in the conditional law of $w_j$.

At each step, the coarse field $e^{\theta_{j-1}}$ lies in $B_r$, after the change of gauge of \cref{ch12:lem:backward-law} if it is given in $D_r$. By the strict domain inclusions established in \cref{ch12:lem:deterministic-prolongation-reserve,ch12:thm:representative-prefixes}, the reconstruction of \cref{ch10:thm:whole-image-gauge} and the reconstruction maps act on the same product. All segments of coarse fields and of coupling-scaled fluctuation inputs used below lie in the larger domains on which the derivative bounds hold. We use \cref{ch13:thm:centered-shell} with $R_{\mathrm{supp}}=r$, together with its change-of-gauge assertion.

\paragraph{Norms.}
Every array is embedded as a componentwise constant function on the half-open cells of its own lattice, each positive link being attached to its starting vertex. We use the heat norm
\begin{equation}\label{ch14:eq:heat-norm}
 \norm f_{-\alpha;T_0}=\sup_{0<s\le T_0}s^{\alpha/2}\norm{P_sf}_\infty,
\end{equation}
where $P_s$ is the continuum heat semigroup acting componentwise; it keeps the constant mode. Since the heat kernel is positive and has unit mass, we have
\begin{equation}\label{ch14:eq:bounded-heat}
 \norm f_{-\alpha;T_0}\le T_0^{\alpha/2}\norm f_\infty\qquad(\alpha>0).
\end{equation}
Discrete heat semigroups carry their mesh as a superscript, as in \cref{ch13:chapter}. For linear and quadratic coefficient arrays, we use the row sums and the volume-normalized column sums of \cref{ch01:def:analytic-norms}: a column of an input on a lattice of mesh $a_{\mathrm{in}}$ is multiplied by $a_{\mathrm{out}}^3/a_{\mathrm{in}}^3$.

\paragraph{Constants.}
We vary the depth at fixed $L$, $T_0$, $H$, and $g_3$, and all limits are taken at these fixed values. Constants have the dependence specified in each result: the lemma on conditional means allows dependence on the upper bound on the coupling, the lemma on column cancellation has the independence stated there, and the theorem on quadratic moments allows dependence on the fixed positive $g_3$ and the fixed $H$. No stronger statement, uniform as $H$ varies or as $g_3\downarrow0$, is claimed; the bounds of \cref{ch12:lem:deterministic-prolongation-reserve} that are uniform in $H$ are not used to improve this dependence. A subscript $J$ indicates a constant that may depend on $J$; such constants are chosen after $J$. The applications of \cref{ch13:lem:heat-and-chaining,ch13:lem:cardinal-columns} use the horizon $T=T_0$, and their constants therefore depend on $T_0/L^2$.

\section{Conditional means along the reconstruction}\label{ch14:sec:means}

The first lemma bounds the reconstructed fields and their conditional mean corrections. The stopping-field contribution to the reconstruction amplitude bound is $r(a_j/H)^{\kapterm}$, where $\kapterm$ is determined by the contraction factor $a_{\mathrm{cap}}=C_{\mathrm{CP}}/B$ of the linearized reconstruction; the fluctuation contribution is bounded by $C_Jg_3\sqrt{a_j}\,p_j$. The conditional mean correction is at most $C_Jg_3^2$ times the size of the coarser field. Indeed, by \cref{ch10:thm:actual-means-and-scores}, the conditional mean of the reconstructed link variables differs from the zero-fluctuation reconstruction by $O(t_j^2\norm{\theta_{j-1}}_\infty)$, and $t_j^2/a_j$ is a fixed multiple of $g_3^2$.

\begin{lemma}\label{ch14:lem:backward-caps-and-means}
Assume the setting of \cref{ch14:sec:setting}, and let the constant $C_{\mathrm{CP}}\ge1$, fixed before $J$, satisfy
\[
 a_{\mathrm{cap}}=C_{\mathrm{CP}}/B,\qquad a_{\mathrm{cap}}\sqrt B<1,\qquad\kapterm=-\frac{\log a_{\mathrm{cap}}}{\log B}>\frac12.
\]
For $1\le j\le k$ set
\[
 A_j=a_j^{-1}\theta_j,\qquad C_j=a_j^{-1}\widehat\Phi_{d_j^k,j}(\theta_{j-1},0),\qquad m_j=\E[A_j\mid\mathcal F_{j-1}]-C_j,
\]
each embedded as a componentwise constant function on the cells of its own lattice. Then, on every sample of the conditioned Wilson measure,
\[
 b_j=\norm{\theta_j}_\infty\le\overline b_j,\qquad\overline b_j=r(a_j/H)^{\kapterm}+C_Jg_3\sqrt{a_j}\,p_j,\qquad\norm{m_j}_\infty\le C_Jg_3^2b_{j-1}.
\]
For $2\le l\le k$ let
\[
 M_l=C_Jg_3^2\bigl\{ra_{\mathrm{cap}}^{l-1}+g_3\sqrt{a_{l-1}}(p_l+1)\bigr\}.
\]
Then $\sum_{j=l}^k\norm{m_j}_\infty\le M_l$, and $M_l\to0$ as $l\to\infty$, uniformly in $k\ge l$. For every $1/2<\alpha<3/2$ and $u\ge0$,
\[
 \sum_{j=l}^k\E\norm{A_j-C_j}_{-\alpha;T_0}\le C_{J,\alpha}g_3a_l^{\alpha-1/2}\sqrt{\log(2+L/a_l)}+C_{\alpha,T_0}M_l,
\]
and, with probability at least $1-Ce^{-u}$, simultaneously for all terms of this finite sum,
\[
 \sum_{j=l}^k\norm{A_j-C_j}_{-\alpha;T_0}\le C_{J,\alpha}g_3a_l^{\alpha-1/2}\bigl[\sqrt{\log(2+L/a_l)}+\sqrt u\bigr]+C_{\alpha,T_0}M_l.
\]
These estimates also hold after integration over the law of the stopping field under the conditioned Wilson measure, uniformly in $k$. The constants depend on the fixed construction parameters, on $L$, $T_0$, $\alpha$, $H$, and on the upper bound on the coupling, but not on $j$, $k$, the depth, or the choice of the representative of the stopping field.
\end{lemma}

The factors $g_3$ and the powers of the mesh are displayed explicitly. The exponent $\kapterm$ is defined by the contraction of the linearized reconstruction and is unrelated to the exponent $\kapflow=1-3\alpflow/2$ of the flow norms. The lemma controls $A_j-C_j$, the deviation of the reconstructed field from its zero-fluctuation reconstruction; the adjacent difference $A_j-A_{j-1}$ contains in addition the last three terms of \eqref{ch14:eq:increment-decomposition}, which are treated in \cref{ch14:sec:signed,ch14:sec:quadratic}.

\begin{proof}
All constants have the dependence allowed in the statement. A constant with subscript $J$ may include the finitely many ratios of couplings within a step and the factor $(1-a_{\mathrm{cap}}\sqrt B)^{-1}$. The proof has four steps: the reconstruction amplitude bound, its summation, the bound on the conditional mean in the gauge of the reconstruction maps, and the bounds on the centered part.

\paragraph{Step 1: the reconstruction amplitude bound.}
The recurrence and its summation in Steps 1 and 2 are also stated in \cref{ch12:lem:backward-law}. We repeat the short computation because the distinction between the constant $C_{\mathrm{rec},J}$ of the recurrence and the constant $C_{\mathrm{cap},J}$ of its solution is used again in the proof of \cref{ch14:thm:quadratic-ancestor-return}. We first explain the choice of $C_{\mathrm{CP}}$. Let $K_1$, $K_2$, and $K_3$ be the bounds, independent of $J$, from \cref{ch12:lem:deterministic-prolongation-reserve} for the linear term, the quadratic term, and the cubic remainder of the normalized zero-fluctuation reconstruction $\overline A_d$ (see \cref{ch14:lem:signed-physical-columns}) on its fixed inner ball. For a coarse field of norm $R\le1$, the value of $\overline A_d$ is at most $K_1R+K_2R^2+K_3R^3\le(K_1+K_2+K_3)R$. We include this sum in $C_{\mathrm{CP}}$; the working ball is chosen inside the inner ball, with $R\le1$, after $J$ is fixed. No constant that depends on all orders of differentiation at fixed $J$ is used to choose $J$. With this choice, the bound
\[
 \norm{\widehat\Phi_{d,j}(\theta,q)}_\infty\le a_{\mathrm{cap}}\norm\theta_\infty+C_{\mathrm{rec},J}\max_s\norm{q_s}_\infty
\]
of \cref{ch12:lem:backward-law} holds with the same $a_{\mathrm{cap}}$; here $C_{\mathrm{rec},J}$ is the constant of that lemma for the coupling-scaled fluctuation inputs. 

Set $\varrho=a_{\mathrm{cap}}\sqrt B<1$. Since $\abs{g_{j,s}w_{j,s}}\le t_jp_j$ on the fluctuation cutoff domain, the last bound implies that, at every point of the support of the conditioned measure,
\[
 b_j\le a_{\mathrm{cap}}b_{j-1}+C_{\mathrm{rec},J}t_jp_j,\qquad b_0\le r.
\]
Let $c_B=M^{(J-1)/2}$, so that $t_j=c_Bg_3\sqrt{a_j}$, and fix
\begin{equation}\label{ch14:eq:cap-constants}
 C_{\mathrm{cap},J}\ge\frac{c_BC_{\mathrm{rec},J}}{1-\varrho}.
\end{equation}
Both $C_{\mathrm{rec},J}$ and $C_{\mathrm{cap},J}$ are chosen after $J$, and neither changes $C_{\mathrm{CP}}$. Since $t_j=t_1B^{-(j-1)/2}$, we have
\begin{equation}\label{ch14:eq:p-affine}
 p_j=p_1+c_{\mathrm w}(j-1),\qquad c_{\mathrm w}=\tfrac12A_{\mathrm w}\log B>0;
\end{equation}
in particular $p_j$ is increasing in $j$. Iterating the recurrence, with no exceptional set, and using $a_{\mathrm{cap}}^{j-i}\sqrt{a_i}=\varrho^{j-i}\sqrt{a_j}$ and $p_i\le p_j$ for $i\le j$, we obtain
\[
 b_j\le ra_{\mathrm{cap}}^j+c_BC_{\mathrm{rec},J}g_3\sqrt{a_j}\sum_{i=1}^j\varrho^{j-i}p_i
 \le ra_{\mathrm{cap}}^j+C_{\mathrm{cap},J}g_3\sqrt{a_j}\,p_j.
\]
Since $a_{\mathrm{cap}}^j=B^{-\kapterm j}=(a_j/H)^{\kapterm}$, this is the asserted bound, with a constant $C_J\ge C_{\mathrm{cap},J}$ that may also dominate the constants of the summation and of the mean below. The condition $\kapterm>1/2$ is equivalent to $\varrho<1$.

\paragraph{Step 2: summation of the amplitude bound.}
Let $q_B=B^{-1/2}$ and $j=l+h$. By the amplitude bound, \eqref{ch14:eq:p-affine}, and $a_{j-1}=a_{l-1}B^{-h}$, we have
\[
 \sum_{j=l}^kb_{j-1}
 \le\frac{ra_{\mathrm{cap}}^{l-1}}{1-a_{\mathrm{cap}}}+C_{\mathrm{cap},J}g_3\sqrt{a_{l-1}}\sum_{h\ge0}q_B^h(p_l+c_{\mathrm w}h)
 \le C_J\bigl\{ra_{\mathrm{cap}}^{l-1}+g_3\sqrt{a_{l-1}}(p_l+1)\bigr\}.
\]
The two series equal $(1-q_B)^{-1}$ and $q_B(1-q_B)^{-2}$. The constant does not depend on $k$ or on the sample.

\paragraph{Step 3: the conditional mean in the gauge of the reconstruction maps.}
Conditionally on $\mathcal F_{j-1}$, the coarse field is fixed, the gauge transformation of the reconstruction maps depends only on it, and the conditional fluctuation kernel is the same as for the reconstruction of \cref{ch10:thm:whole-image-gauge}. Equality of the kernels does not by itself identify the two reconstructed means, because the observable changes; we show that the bound on the mean survives the change of gauge. We write $\phi=\theta_{j-1}$ and $R=\norm\phi_\infty$. Let $X_e(w)$ be the logarithm of the reconstruction of \cref{ch10:thm:whole-image-gauge} on one output link $e$, and let $x_e=X_e(0)$. By the bound on the reconstructed mean in \cref{ch10:thm:actual-means-and-scores}, we have
\[
 \abs{\E X_e-x_e}\le C_Jg_j^2R,
\]
where we used $t_j^2=c_B^2g_j^2$. The variance of $X_e$ is at most $C_Jg_j^2$, by \cref{ch13:thm:centered-shell}, applied under the conditional fluctuation kernel at the coarse field $e^\phi\in B_r$, with fine mesh $a_j$ and horizon $T_0$. Indeed, we apply its variance bound to the test function $a_j^{-3}\one_{Q_e}$ in one color component of $X/a_j$, whose squared $L^2$ norm is $a_j^{-3}$, and multiply by $a_j^2$. Summing over the three colors, we obtain
\begin{equation}\label{ch14:eq:original-second-moment}
 \E\abs{X_e-x_e}^2\le C_Jg_j^2+C_Jg_j^4R^2\le C_J'g_j^2.
\end{equation}
The reconstruction of \cref{ch10:thm:whole-image-gauge} satisfies the hypotheses of both theorems used here.

At the fixed coarse field, the reconstruction maps differ from this reconstruction by a gauge transformation given by matrices $h_x$ at the vertices. On the link $e=(x,y)$, the logarithm of the reconstruction map is $H_e(X_e)$, where $H_e(X)=\Log(h_xe^Xh_y^{-1})$. By the bounds of \cref{ch10:thm:whole-image-gauge,ch12:thm:replacement-family}, the two zero-fluctuation reconstructions $x_e$ and $\widehat x_e=H_e(x_e)$ have norm at most $CR$. The identity $e^{\widehat x_e}=h_xe^{x_e}h_y^{-1}$ therefore implies $\abs{\Log(h_y^{-1}h_x)}\le CR$. Let $z_e=\Log(h_y^{-1}h_x)$. By conjugation, we have
\[
 H_e(X)=\Ad(h_x)\Log(e^Xe^{z_e}).
\]
On the fixed small ball of logarithms, the analytic series of the exponential and the logarithm give bounded derivatives of $(X,z)\mapsto\Log(e^Xe^z)$ through order three. More precisely, we enlarge this ball within the fixed chart so that $\norm{e^Xe^z-I}\le q<1$ in matrix norm. Termwise differentiation of $\Log(I+Z)=\sum_{m\ge1}(-1)^{m+1}Z^m/m$ is justified by the convergent majorants $\sum_{m\ge r}m^{r-1}q^{m-r}$, $r=1,2,3$, and the derivatives of the exponential are bounded by its absolutely convergent series on this ball. These bounds depend only on the chart. At $z=0$, the function equals $X$, and its second derivative in $X$ vanishes there. Integrating the mixed derivative $D_zD_X^2$ along the segment $sz_e$, $0\le s\le1$, we therefore find $\norm{D_X^2H_e}\le CR$, while $\norm{D_XH_e}\le C$. All segments used here lie in the fixed chart. Since the matrices $\Ad(h_x)$ are orthogonal, these constants do not depend on the matrices $h_x$ or on the number of vertices. By Taylor's formula at $x_e$, the bound on the mean, and \eqref{ch14:eq:original-second-moment}, we obtain
\[
 \abs{\E H_e(X_e)-H_e(x_e)}\le C\abs{\E(X_e-x_e)}+CR\,\E\abs{X_e-x_e}^2\le C_Jg_j^2R.
\]
This is the bound on the mean for the reconstruction maps, on the same product of balls. Dividing by the fine mesh, we obtain
\[
 \norm{m_j}_\infty\le C_J\frac{t_j^2}{a_j}b_{j-1}\le C_Jg_3^2b_{j-1}.
\]
Together with Step 2, this implies $\sum_{j=l}^k\norm{m_j}_\infty\le M_l$, after one enlargement of $C_J$. Both terms of $M_l$ tend to zero as $l\to\infty$, because $a_{\mathrm{cap}}<1$ and $p_l$ grows linearly in $l$ while $\sqrt{a_{l-1}}$ decays geometrically.

\paragraph{Step 4: the centered part.}
We define
\[
 F_j=A_j-\E[A_j\mid\mathcal F_{j-1}].
\]
\Cref{ch13:thm:centered-shell} applies to $F_j$ conditionally on $\mathcal F_{j-1}$. Indeed, the fine mesh is $a_j$, the side of the coarse lattice of the step is $n_0B^{j-1}\ge4$, the depth is $d_j^k$, the largest coupling is $t_j\le g_*$, the coarse field is in $B_r$, the horizon is $T_0$, and the reconstruction maps satisfy the hypotheses of that theorem on the whole product. The same holds for the reconstruction of \cref{ch10:thm:whole-image-gauge}, which we used in Step 3. Let $\beta=\alpha-1/2>0$ and $\ell_j=\log(2+L/a_j)$. Taking the conditional expectation and then the expectation, we obtain
\[
 \E\norm{F_j}_{-\alpha;T_0}\le C_{J,\alpha}g_3a_j^\beta\sqrt{\ell_j}.
\]
For $h\ge0$, we have $\ell_{l+h}\le\ell_l+h\log B$ and $a_{l+h}=a_lB^{-h}$. It follows that
\begin{equation}\label{ch14:eq:geometric-entropy}
 \sum_{h\ge0}a_{l+h}^\beta\sqrt{\ell_{l+h}}
 \le a_l^\beta\Bigl\{\frac{\sqrt{\ell_l}}{1-B^{-\beta}}+\sqrt{\log B}\sum_{h\ge0}B^{-\beta h}\sqrt h\Bigr\}
 \le C_{B,\beta}a_l^\beta\sqrt{\ell_l}.
\end{equation}
The last series converges, for instance because $\sqrt h\le h+1$, and $\ell_l\ge\log2$ absorbs its contribution. Since $A_j-C_j=F_j+m_j$, the bound on the expectation follows from \eqref{ch14:eq:bounded-heat} applied to $m_j$ and from Step 3.

For the bound in probability, we use the tail bound of \cref{ch13:thm:centered-shell} at level $j=l+h$ with $u+2h\log2$ in place of $u$. Integrating the conditional probability does not change its upper bound. By a union bound over the finitely many levels, the probability of failure is at most
\[
 Ce^{-u}\sum_{h\ge0}4^{-h}\le C'e^{-u}.
\]
On the complementary event, we sum the thresholds, using $\sqrt{x+y}\le\sqrt x+\sqrt y$, \eqref{ch14:eq:geometric-entropy}, and the convergent series $\sum_hB^{-\beta h}$ and $\sum_hB^{-\beta h}\sqrt h$, and we add the deterministic sum of the means by \eqref{ch14:eq:bounded-heat}. This proves the bound in probability. No independence between levels has been used.

Since all conditional estimates are uniform in the coarse fields that occur and in the stopping field, the same inequalities hold after integration over the law of the stopping field under the conditioned Wilson measure. The argument concerns $A_j-C_j$ throughout, and $C_j$ is never replaced by $A_{j-1}$.
\end{proof}

\section{Column cancellation and heat-kernel estimates}\label{ch14:sec:signed}

The normalization of an input column is essential here. For an operator $K$ from arrays on a lattice of mesh $A$ to componentwise functions on the torus, the volume-normalized absolute column of an input $x$ is $A^{-3}\int\abs{K_x(y)}\,dy$. A zero integral of a column is a separate identity, which no absolute bound implies. On a lattice of mesh $\delta$, the integral is replaced by $\delta^3\sum_y$, still divided by the input volume $A^3$.

The lemma below concerns the difference $D_d$ between the linearized zero-fluctuation reconstruction and the cell embedding on the coarse lattice. Each column of $D_d$ is concentrated within distance of order $A$ of its input, and it has zero integral. At times $t\gg A^2$, the heat kernel is nearly constant on the scale $A$, and the zero integral allows us to replace $p_t(z-y)$ by $p_t(z-y)-p_t(z-x)$, which costs the distance from $y$ to $x$ times the gradient of the kernel. This yields the factor $A/\sqrt t$.

\begin{lemma}\label{ch14:lem:signed-physical-columns}
Assume the setting of \cref{ch14:sec:setting}. Let $d\ge0$, let $A>0$ be a mesh with $L/A\in\N$ and $L/A\ge4$, let $a=A/B$, and let $\widehat\Phi_d$ be the logarithmic form of the reconstruction map of one renormalization group step at depth $d$, from the lattice of mesh $A$ to the lattice of mesh $a$. In units of the coarse lattice, its zero-fluctuation reconstruction has the expansion
\[
 \widehat\Phi_d(\theta,0)=B^{-1}\overline A_d(\theta),\qquad\overline A_d(\theta)=\overline P_d\theta+\overline Q_d(\theta,\theta)+\overline R_d(\theta),
\]
where $\overline P_d=D\overline A_d(0)$, $\overline Q_d=D^2\overline A_d(0)/2$, and $\norm{\overline R_d(\theta)}_\infty\le C_{\mathrm{CP}}\norm\theta_\infty^3$ on a fixed inner ball of coarse fields. The first three derivatives in $\theta$, together with their first spatial differences $\nabla_{1/B}$ in units of the coarse lattice, have weighted row sums and volume-normalized column sums bounded by one constant, uniformly in $d$, the lattice side, and $J$, as in \cref{ch12:lem:deterministic-prolongation-reserve}. Here $\nabla_{1/B}=A\nabla_a$, where $\nabla_a$ denotes the first differences at the physical mesh $a$; thus a physical first difference of $\overline A_d$ carries the factor $A^{-1}$, and a physical first difference of the potential $A^{-1}\overline A_d$ carries the factor $A^{-2}$. The cubic remainder and these finitely many bounds are included in $C_{\mathrm{CP}}$ before $J$ is fixed. In physical units the zero-fluctuation reconstruction is $A^{-1}\overline A_d(\theta)$.

Let $E_Av$ be the cell embedding on the lattice of mesh $A$ of a coarse lattice potential $v$, let $P_dv$ be the cell embedding on the lattice of mesh $a$ of the array $A^{-1}\overline P_d(Av)$, and let $D_d=P_d-E_A$. For each input link $x$ and input color $b$, and each output direction and color $(i,c)$, the corresponding column of $D_d$ satisfies
\[
 \int_{\Torus}(D_d)_{(i,c),(x,b)}(y)\,dy=0,\qquad A^{-3}\int_{\Torus}\dist(y,x)\bigl\lvert(D_d)_{(i,c),(x,b)}(y)\bigr\rvert\,dy\le C_{\mathrm{CP}}A.
\]
Its volume-normalized column sums and its row sums are bounded by $C_{\mathrm{CP}}$, with exponential localization at the scale $A$. In particular, for $0<t\le T_0$,
\[
 \norm{P_tD_dv}_\infty\le C\min\{1,A/\sqrt t\}\norm v_\infty,\qquad\norm{D_dv}_{-\alpha;T_0}\le C_\alpha A^\alpha\norm v_\infty\quad(0<\alpha<1).
\]
Let $m\in\N$ and $\delta=a/m$. Let $P_d^{\mathrm{arr}}v$ be the array on the lattice of mesh $a$ whose cell embedding is $P_dv$, and let
\[
 D_d^\delta=I_{a\to\delta}P_d^{\mathrm{arr}}-I_{A\to\delta},
\]
where the two cardinal interpolations of \cref{ch13:lem:cardinal-columns} start from the lattices of meshes $a$ and $A$, respectively. Then every componentwise column of $D_d^\delta$ has zero integral $\delta^3\sum_y$, its absolute column sums and first moments satisfy the same bounds, and for every integer $q\ge0$, every product $D_\delta^q$ of $q$ physical forward or backward differences, and every $0\le t\le T_0$,
\[
 \norm{D_\delta^qP_t^\delta D_d^\delta v}_\infty\le C_{J,q}(t+A^2)^{-q/2}\min\bigl\{1,A/\sqrt{t+A^2}\bigr\}\norm v_\infty.
\]
In the continuum estimates, $C$ depends only on the column and moment bounds, on $\alpha$ where it appears, and on $T_0/L^2$. In the estimates on the lattice of mesh $\delta$ the constant may also depend on $J$, $q$, and the function $b$ of the cardinal interpolation. No constant depends on $\delta/a$, the lattice side, the depth, the coupling, or the coarse field. The zero integral is a linear identity for each input column; no zero integral of $\overline Q_d$ is asserted.
\end{lemma}

\begin{proof}
The proof has five steps: the expansion, the zero integrals, the localization, the continuum heat estimates, and the estimates on the finer lattice.

\paragraph{Step 1: the expansion and its units.}
By \cref{ch12:lem:deterministic-prolongation-reserve}, the reconstruction maps keep the bounds on the first three derivatives of the zero-fluctuation reconstruction in the coarse field and on its first spatial difference in units of the coarse lattice, the volume-normalized columns, and the zero-integral identity of \cref{ch12:thm:replacement-family}. Their common domain contains a fixed inner ball of coarse fields. Taylor's formula on this ball gives
\[
 \begin{gathered}
 \overline A_d(\theta)=D\overline A_d(0)\theta+\tfrac12D^2\overline A_d(0)[\theta,\theta]+\overline R_d(\theta),\\
 \overline R_d(\theta)=\tfrac12\int_0^1(1-s)^2D^3\overline A_d(s\theta)[\theta,\theta,\theta]\,ds.
 \end{gathered}
\]
The value at $\theta=0$ is zero by global color equivariance. The bound on the third derivative implies $\norm{\overline R_d(\theta)}_\infty\le(1/6)\sup\norm{D^3\overline A_d}\,\norm\theta_\infty^3$. We enlarge $C_{\mathrm{CP}}$ to include this bound, the finitely many derivative bounds, and their first spatial differences in units of the coarse lattice, before $J$ is fixed. This uses the finitely many constants that are uniform in $J$, not the bounds on all derivatives at fixed $J$. The conversion to physical units is
\begin{equation}\label{ch14:eq:center-units}
 a^{-1}\widehat\Phi_d(\theta,0)=(A/B)^{-1}B^{-1}\overline A_d(\theta)=A^{-1}\overline A_d(\theta).
\end{equation}

\paragraph{Step 2: zero integrals.}
We first restate the zero-integral identity of \cref{ch12:thm:replacement-family} in these units. The reconstruction satisfies the averaging constraint exactly. Differentiating this identity at $\theta=0$, we find that the linearized reconstruction preserves the componentwise physical integral of each input column. For each physical input basis vector, at an input link $x$, with a direction and a color, this gives
\[
 \int_{\Torus}(P_d)_{(i,c),(x,b)}(y)\,dy=\int_{\Torus}(E_A)_{(i,c),(x,b)}(y)\,dy.
\]
The direction of $x$ is suppressed in the notation, as in the statement. The matching component on the right has integral $A^3$, and the other components have integral zero. Subtracting, we obtain the zero integral of each column of $D_d$. This is the exact identity of \cref{ch12:thm:replacement-family}, not a consequence of its absolute bounds. In that identity, the periodic gradient part of the differentiated reconstruction has zero sum, including across the seams of the torus; hence no constant mode has been discarded.

\paragraph{Step 3: localization and the first moment.}
After a fixed reduction of the exponent of the spatial weight, there exist constants $C,c>0$, independent of $J$, such that for every column and every integer $h\ge0$,
\begin{equation}\label{ch14:eq:column-shell}
 A^{-3}\int_{hA\le\dist(y,x)<(h+1)A}\abs{(D_d)_x(y)}\,dy\le Ce^{-ch},
\end{equation}
and the row sums are bounded. These statements follow by subtracting the coarse cell column from the volume-normalized column of $P_d$, weighted exponentially; moving a point of a cell to the vertex of the cell costs at most $\sqrt3A$. The same argument, applied to each weighted row, yields the localization of the rows. Multiplying \eqref{ch14:eq:column-shell} by $(h+1)A$ and summing over $h$, we obtain
\[
 A^{-3}\int\dist(y,x)\abs{(D_d)_x(y)}\,dy\le CA\sum_{h\ge0}(h+1)e^{-ch}\le C'A.
\]
All sums over directions and colors are finite and included in $C$.

\paragraph{Step 4: heat-kernel cancellation in the continuum.}
In this step, all heat estimates use the horizon $T=T_0$ in \cref{ch13:lem:heat-and-chaining}. To pass from the moments of individual columns to a bound on the supremum over outputs, we use the following elementary inequality. For a smooth periodic function $f$ and a cube $Q$ of side $A$,
\begin{equation}\label{ch14:eq:cube-sup}
 A^3\sup_Q\abs f\le C\sum_{S\subset\{1,2,3\}}A^{\abs S}\int_Q\abs{\partial_Sf(y)}\,dy.
\end{equation}
Indeed, on an interval of length $A$, we have $\abs{f(x)}\le A^{-1}\int\abs f+\int\abs{f'}$, and applying this successively in the three variables yields \eqref{ch14:eq:cube-sup}; the endpoints are handled by continuity. If $Q$ is enlarged to the set of points within distance $(h+2)A$, the enlargement is covered by at most $C(h+3)^3$ cubes of side $A$. Summing the resulting bounds on the supremum over the lattice of inputs costs the same overlap factor $C(h+3)^3$.

Let $A^2\le t\le T_0$ and $f_z(y)=p_t(z-y)$. Since every column has zero integral, we have
\[
 (P_tD_dv)(z)=\sum_xv_x\int\bigl[p_t(z-y)-p_t(z-x)\bigr](D_d)_x(y)\,dy.
\]
On the $h$th annulus around $x$, the bracket is at most $(h+1)A$ times the supremum of $\abs{\nabla p_t(z-\cdot)}$ on the enlarged cube around $x$; a shortest periodic path proves this also when the annulus crosses a seam. Using \eqref{ch14:eq:column-shell}, then \eqref{ch14:eq:cube-sup} with the overlap bound, and finally the $L^1$ bounds on the derivatives of the heat kernel in \cref{ch13:lem:heat-and-chaining}, we obtain
\[
 \norm{P_tD_dv}_\infty\le CA\norm v_\infty\sum_{h\ge0}e^{-ch}(h+3)^4\sum_{r=0}^3A^r\norm{\nabla^{r+1}p_t}_1\le C\frac A{\sqrt t}\norm v_\infty\qquad(A^2\le t\le T_0),
\]
where in the last inequality we used that each factor $(A/\sqrt t)^r$ is at most one. For $0<t\le\min\{A^2,T_0\}$, the contraction of the heat semigroup and the bound on the row sums imply $\norm{P_tD_dv}_\infty\le C\norm v_\infty$. This proves the value estimate, with constants depending only on the column bounds and on $T_0/L^2$. When $0<\alpha<1$, multiplying by $t^{\alpha/2}$ gives at most $CA^\alpha$ in both ranges of time; if $A^2>T_0$, the bound for small times suffices. This proves the bound on the heat norm.

\paragraph{Step 5: the finer lattice.}
We now keep the lattices of meshes $a$ and $A$ distinct. Let $m\in\N$ and $\delta=a/m$, so that
\begin{equation}\label{ch14:eq:refinement-ratios}
 \frac a\delta=m\in\N,\qquad\frac A\delta=Bm\in\N.
\end{equation}
Since $L/A\in\N$ and $L/A\ge4$, the numbers $L/a=B(L/A)$ and $L/\delta=Bm(L/A)$ are integers at least four. We therefore apply \cref{ch13:lem:cardinal-columns} to the pairs of meshes $(a,\delta)$ and $(A,\delta)$, with integer ratios $m$ and $Bm$, the same function $b$, and the horizon $T=T_0$. For a fixed integer $q\ge0$, $D_\delta^q$ denotes any ordered product of $q$ physical forward or backward differences, and the identity if $q=0$; these are exactly the products allowed in both lemmas of \cref{ch13:chapter}. Let $p_{z,x}$ be the coefficients of the array $P_d^{\mathrm{arr}}$. By the mass identity of \cref{ch13:lem:cardinal-columns}, we have, for each componentwise column,
\begin{equation}\label{ch14:eq:cardinal-signed-mass}
 \begin{split}
 \delta^3\sum_y(D_d^\delta)_{y,x}
 &=\delta^3\sum_y\sum_{z\in\Lat{L/a}}\zeta_3((y-z)/a)p_{z,x}-\delta^3\sum_y\zeta_3((y-x)/A)\Id\\
 &=a^3\sum_zp_{z,x}-A^3\Id=0,
 \end{split}
\end{equation}
where the last equality is the discrete form of the zero integral in Step 2. The identity matrices in this formula involve the matching direction and color only. Since the function $b$ is nonnegative, the mass and moment bounds of the interpolation imply
\[
 \begin{split}
 \delta^3\sum_y\Bigl\lvert\sum_z\zeta_3((y-z)/a)p_{z,x}\Bigr\rvert&\le a^3\sum_z\abs{p_{z,x}},\\
 \delta^3\sum_y\dist(y,x)\Bigl\lvert\sum_z\zeta_3((y-z)/a)p_{z,x}\Bigr\rvert&\le a^3\sum_z\bigl(\dist(z,x)+Ca\bigr)\abs{p_{z,x}}.
 \end{split}
\]
The interpolated coarse column has mass $A^3$ and first moment at most $CA^4$. Both terms of $D_d^\delta$ therefore satisfy the bounds of the columns before interpolation, with a fixed enlargement of the constants. Since their supports are displaced by at most $Ca\le CA$, \eqref{ch14:eq:column-shell} holds with the integral replaced by $\delta^3\sum_y$. The row sums are bounded as well, by the positive partition of unity of each interpolation.

The argument of Step 4 applies on the lattice of mesh $\delta$, and we only describe its new ingredient. On an interval of $m_A=A/\delta=Bm$ lattice points, averaging the inequality
\[
 \abs{f(r)}\le\abs{f(s)}+\delta\sum_{u\text{ between }r,s}\abs{D_\delta f(u)}
\]
over $s$ yields the one-dimensional analog of \eqref{ch14:eq:cube-sup}. Taking the product over the coordinates, with an enlargement by one lattice step to include the differences at the boundary, we obtain \eqref{ch14:eq:cube-sup} with sums in place of integrals and products of forward differences in place of derivatives. The overlap bound for the enlarged cubes is unchanged. To subtract the kernel, we use \eqref{ch14:eq:cardinal-signed-mass} and a shortest lattice path, whose length is at most a fixed multiple of the torus distance. Then the $L^1$ bounds on discrete differences of the heat kernel in \cref{ch13:lem:heat-and-chaining} imply
\begin{equation}\label{ch14:eq:discrete-defect-value}
 \norm{P_t^\delta D_d^\delta v}_\infty\le C\min\{1,A/\sqrt t\}\norm v_\infty\qquad(0<t\le T_0).
\end{equation}
No constant here depends on $a/\delta$.

For $0\le t\le\min\{2A^2,T_0\}$, we commute the product of $q$ differences with the heat semigroup and use the derivative bounds of the two interpolations:
\[
 \norm{D_\delta^qP_t^\delta D_d^\delta v}_\infty
 \le\norm{D_\delta^qI_{a\to\delta}P_d^{\mathrm{arr}}v}_\infty+\norm{D_\delta^qI_{A\to\delta}v}_\infty
 \le C_q(a^{-q}+A^{-q})\norm v_\infty\le C_{J,q}A^{-q}\norm v_\infty.
\]
For $2A^2\le t\le T_0$, we split the heat semigroup into two equal halves, use the $L^1$ bound on the differentiated heat kernel for the first half and \eqref{ch14:eq:discrete-defect-value} for the second:
\[
 \norm{D_\delta^qP_t^\delta D_d^\delta v}_\infty\le C_q(t/2+\delta^2)^{-q/2}\frac A{\sqrt{t/2}}\norm v_\infty\le C_qAt^{-(q+1)/2}\norm v_\infty.
\]
Up to fixed constants, the two ranges of time give the asserted bound with $t+A^2$; they include $t=0$ and every fixed order of differences. Only the estimate for small times introduces the factor $B^q$, through $a^{-q}=B^qA^{-q}$, and this factor is allowed in $C_{J,q}$. All the cancellations used in this proof are identities for the columns of linear operators; neither Taylor's formula nor the heat-kernel cancellation yields a zero integral of $\overline Q_d$.
\end{proof}

\section{Conditional quadratic moments}\label{ch14:sec:quadratic}

\paragraph{Quadratic maps and their norms.}
Let $Q$ be a quadratic map from arrays on a lattice of mesh $a_{\mathrm{in}}$ to arrays on a lattice of mesh $a_{\mathrm{out}}\le a_{\mathrm{in}}$. We write a scalar output component as
\[
 Q_y(v,v)=\sum_{x,z}T_{y;x,z}v_xv_z,
\]
where the indices include directions and colors. We use the symmetric polarization: replacing $T_{y;x,z}$ by $(T_{y;x,z}+T_{y;z,x})/2$ leaves $Q(v,v)$ unchanged and does not increase any of the coefficient sums below. We write $Q(u,v)$ for the corresponding symmetric bilinear map, and $Q:\mathcal M$ for the array with components $\sum_{x,z}T_{y;x,z}\mathcal M_{xz}$, for a matrix $\mathcal M$ indexed by the inputs. Let $\Tree(y,x,z)$ be the length of a shortest connected network in $\Torus$ containing the locations of $y$, $x$, and $z$, with internal vertices allowed, as in \cref{ch01:def:analytic-norms}. We fix a spatial exponent $\asp'>0$, smaller than the exponents of the coefficient bounds of \cref{ch12:chapter} by a fixed amount, and set
\begin{equation}\label{ch14:eq:weighted-quadratic-norm}
 \norm T_{a_{\mathrm{in}},a_{\mathrm{out}}}=\max\Bigl\{\sup_y\sum_{x,z}e^{\asp'\Tree(y,x,z)/a_{\mathrm{in}}}\abs{T_{y;x,z}},\ \sup_x\frac{a_{\mathrm{out}}^3}{a_{\mathrm{in}}^3}\sum_{y,z}e^{\asp'\Tree(y,x,z)/a_{\mathrm{in}}}\abs{T_{y;x,z}}\Bigr\}.
\end{equation}
By symmetry, the column with the second input fixed has the same bound. In particular, if $a_{\mathrm{in}}=a_i$, $a_{\mathrm{out}}=\delta$, and $\norm T_{a_i,\delta}\le K_Q$, then
\begin{equation}\label{ch14:eq:quadratic-norm}
 \sup_y\sum_{x,z}\abs{T_{y;x,z}}\le K_Q,\qquad\sup_x\frac{\delta^3}{a_i^3}\sum_{y,z}\abs{T_{y;x,z}}\le K_Q.
\end{equation}
The same exponent $\asp'$ is used at every level below; it is not reduced after each composition. At the coarsest level, we set $\overline b_0=r$. For $i\ge1$, $\overline b_i$ always denotes the reconstruction amplitude bound of \cref{ch14:lem:backward-caps-and-means}, never a random bound left outside a conditional expectation.

\paragraph{Why a conditional decomposition is needed.}
For $j\ge2$, by the reconstruction amplitude bound, the quadratic term $a_{j-1}^{-1}\overline Q(\theta_{j-1},\theta_{j-1})$ in \eqref{ch14:eq:increment-decomposition} is at most $Ca_{j-1}^{-1}\overline b_{j-1}^2$, which is in general of order $g_3^2p_{j-1}^2$ and is not summable in $j$. We therefore regard $\theta_{j-1}$ as reconstructed from $\theta_{j-2}$, condition on $\mathcal F_{j-2}$, and separate a conditionally centered part, which concentrates, from a conditional expectation error, which is smaller. The conditional expectation is close to the quadratic term evaluated on the linearized reconstruction $L_{j-1}\theta_{j-2}$, and repeating the argument pulls the quadratic coefficient back through the linearized reconstructions. Each pullback contracts the coefficient by the factor $a_{\mathrm{cap}}^2$, one factor for each argument, while dividing by the coarser mesh costs a factor $B$; the resulting ratio $\qanc=Ba_{\mathrm{cap}}^2<1$ makes the final convolution summable. Volume-normalized columns are essential here: with counting column sums, each pullback would lose an additional factor $B^3$.

Three errors remain in the conditional expectation, and they explain the three terms in the first part of \cref{ch14:thm:quadratic-ancestor-return} below. Replacing the nonlinear zero-fluctuation reconstruction by its linearization costs $a_{\mathrm{cap}}^2\norm\phi_\infty^3$. The conditional means of the fluctuation variables and the change of their covariance with the coarse field cost $g_i^2\norm\phi_\infty$. Here we use that the conditional kernel at the zero coarse field is invariant under constant color rotations, although it is spatially correlated; an expected adjoint vector built from it therefore vanishes. The terms of higher order in the fluctuations cost $g_i^3p_i^3$.

The deterministic part of the pullback is the following lemma.

\begin{lemma}\label{ch14:lem:quadratic-pullback}
Let $1\le i\le j$, let $T$ be the symmetric coefficient array of a quadratic map from arrays on the lattice of mesh $a_i$ to arrays on the lattice of mesh $a_j$, and let $L$ be the coefficient array of a linear map from arrays on the lattice of mesh $a_{i-1}=Ba_i$ to arrays on the lattice of mesh $a_i$. Suppose that $\norm T_{a_i,a_j}\le K$ and that
\begin{equation}\label{ch14:eq:tangent-rows-columns}
 \sup_x\sum_ue^{\asp'\dist(x,u)/a_{i-1}}\abs{L_{x,u}}\le a_{\mathrm{cap}},\qquad\sup_u\frac{a_i^3}{a_{i-1}^3}\sum_xe^{\asp'\dist(x,u)/a_{i-1}}\abs{L_{x,u}}\le a_{\mathrm{cap}}.
\end{equation}
Then the coefficient array
\[
 T'_{y;u,v}=\sum_{x,z}T_{y;x,z}L_{x,u}L_{z,v}
\]
of the map $v\mapsto Q(Lv,Lv)$ is symmetric and satisfies $\norm{T'}_{a_{i-1},a_j}\le Ka_{\mathrm{cap}}^2$. If the quadratic map is equivariant under the adjoint action of constant color rotations and $L$ commutes with them, then the map $v\mapsto Q(Lv,Lv)$ is equivariant.
\end{lemma}

\begin{proof}
We first prove the unweighted bounds. The output row of $T'$ is at most $Ka_{\mathrm{cap}}^2$, by the row bounds of $T$ and of $L$. For the column with the first input $u$ fixed, we use the volume-normalized column of $L$ in the first argument and its row in the second:
\[
 \frac{a_j^3}{a_{i-1}^3}\sum_{y,v}\abs{T'_{y;u,v}}
 \le\frac{a_i^3}{a_{i-1}^3}\sum_x\abs{L_{x,u}}\,a_{\mathrm{cap}}\Bigl(\frac{a_j^3}{a_i^3}\sum_{y,z}\abs{T_{y;x,z}}\Bigr)
 \le Ka_{\mathrm{cap}}^2.
\]
Fixing the second input instead repeats this calculation with $x,z$ and $u,v$ interchanged. The factor $a_i^3/a_{i-1}^3=B^{-3}$ cancels the factor $B^3$ in the counting column sum of $L$.

The weights cause no additional loss. A shortest network for $(y,x,z)$, joined with shortest paths from $x$ to $u$ and from $z$ to $v$, is a connected network containing $(y,u,v)$, which implies $\Tree(y,u,v)\le\Tree(y,x,z)+\dist(x,u)+\dist(z,v)$; periodic shortest networks satisfy the same inequality. In the units $a_{i-1}>a_i$ of the new input mesh, the weight of the old network is at most its weight in the units $a_i$, and the weights in \eqref{ch14:eq:tangent-rows-columns} are already in the units $a_{i-1}$. The new weight is therefore at most the product of the three old weights, and the preceding calculation applies with these weights inserted and the same exponent. Symmetry of $T'$ follows from that of $T$. Equivariance is preserved because $L$ commutes with the color rotations.
\end{proof}

\begin{theorem}\label{ch14:thm:quadratic-ancestor-return}
Assume the setting of \cref{ch14:sec:setting} and fix $k\ge1$. Let $\overline b_i$, $1\le i\le k$, be the reconstruction amplitude bounds of \cref{ch14:lem:backward-caps-and-means}, and $\overline b_0=r$. Let $a_{\mathrm{cap}}=C_{\mathrm{CP}}/B$ and $\qanc=Ba_{\mathrm{cap}}^2=C_{\mathrm{CP}}^2/B<1$, and for $1\le i\le k$ let $L_i=D_\theta\widehat\Phi_{d_i^k,i}(0,0)$ be the linearized zero-fluctuation reconstruction of the $i$th step.

Let $1\le i\le k$, and let $Q$ be a deterministic quadratic map from arrays on the lattice of mesh $a_i$ to $\su$-valued arrays on a lattice of mesh $\delta\le a_i$, which is equivariant under the adjoint action of constant color rotations and satisfies $\norm T_{a_i,\delta}\le K_Q$ for its coefficient array $T$, with all input and output directions and colors included. Then the following statements hold.
\begin{enumerate}
\item For every coarse field $e^\phi\in B_r$, under the conditional fluctuation kernel of the $i$th step at the coarse field $e^\phi$,
\[
 \bigl\lVert\E[Q(\theta_i,\theta_i)\mid\theta_{i-1}=\phi]-Q(L_i\phi,L_i\phi)\bigr\rVert_\infty
 \le C_JK_Q\bigl\{a_{\mathrm{cap}}^2\norm\phi_\infty^3+g_i^2\norm\phi_\infty+g_i^3p_i^3\bigr\}.
\]
Let $\mu_\phi=\E_\phi w$ and $\Sigma_\phi=\Cov_\phi(w,w)$ be the conditional mean and covariance of the fluctuation variables under this kernel (in the notation of \cref{ch23:chapter}, $\mu_\phi=m_\phi$), and let $E_\phi$ be the derivative at $w=\mu_\phi$ of the map $w\mapsto g_i^{-1}\widehat\Phi_{d_i^k,i}(\phi,(g_{i,s}w_s)_s)$, which includes the ratios $g_{i,s}/g_i$. Then the expectation of $Q(g_iE_\phi(w-\mu_\phi),g_iE_\phi(w-\mu_\phi))$ equals $g_i^2Q:(E_\phi\Sigma_\phi E_\phi^{\mathsf T})$. This contraction vanishes at $\phi=0$, and
\[
 \norm{E_\phi\Sigma_\phi E_\phi^{\mathsf T}-E_0\Sigma_0E_0^{\mathsf T}}_{\mathrm{op}}\le C_J\norm\phi_\infty.
\]
The first estimate accounts for the value of the reconstruction at $\mu_\phi$, its displacement from the zero-fluctuation reconstruction, and all remainders of the means and of the covariance; it does not assume that the conditional expectation of a quadratic term vanishes when $\phi\ne0$.
\item For every $0<\alpha<3/2$, on every sample,
\[
 \E\Bigl[\bigl\lVert Q(\theta_i,\theta_i)-\E[Q(\theta_i,\theta_i)\mid\mathcal F_{i-1}]\bigr\rVert_{-\alpha;T_0}\Bigm|\mathcal F_{i-1}\Bigr]
 \le C_{J,\alpha}K_Qg_i\overline b_ia_i^\alpha\sqrt{\log(2+L/a_i)}.
\]
\item For $2\le j\le k$ define quadratic maps $Q_{j,i}$ from arrays on the lattice of mesh $a_i$ to arrays on the lattice of mesh $a_j$, for $0\le i\le j-1$, by
\[
 Q_{j,j-1}=\overline Q_{d_j^k,j},\qquad Q_{j,i-1}(v,v)=Q_{j,i}(L_iv,L_iv)\quad(1\le i<j),
\]
where $\overline Q_{d_j^k,j}$ is the quadratic part of the zero-fluctuation reconstruction of the $j$th step, as in \cref{ch14:lem:signed-physical-columns}. These maps are equivariant. Let $T_{j,i}$ be the symmetric coefficient array of $Q_{j,i}$, and set $K_{j,i}=\norm{T_{j,i}}_{a_i,a_j}$. Then
\[
 K_{j,i}\le Ca_{\mathrm{cap}}^{2(j-1-i)},\qquad\frac{K_{j,i}}{a_{j-1}}\le C\frac{\qanc^{j-1-i}}{a_i}.
\]
\item For $2\le j\le k$ and $1\le i\le j-1$ put
\[
 \begin{split}
 Z_{j,i}&=Q_{j,i}(\theta_i,\theta_i)-\E[Q_{j,i}(\theta_i,\theta_i)\mid\mathcal F_{i-1}],\\
 e_{j,i}&=\E[Q_{j,i}(\theta_i,\theta_i)\mid\mathcal F_{i-1}]-Q_{j,i-1}(\theta_{i-1},\theta_{i-1}).
 \end{split}
\]
Then, on every sample,
\[
 \overline Q_{d_j^k,j}(\theta_{j-1},\theta_{j-1})=Q_{j,0}(\theta_0,\theta_0)+\sum_{i=1}^{j-1}(Z_{j,i}+e_{j,i}).
\]
Further, with
\[
 v_i=g_i\overline b_ia_i^{\alpha-1},\qquad u_i=\frac{\qanc\overline b_{i-1}^{\,3}}{a_{i-1}}+g_3^2\overline b_{i-1}+g_3^3\sqrt{a_i}\,p_i^3,
\]
we have, for $0<\alpha<3/2$,
\[
 \E\bigl\lVert a_{j-1}^{-1}\overline Q_{d_j^k,j}(\theta_{j-1},\theta_{j-1})\bigr\rVert_{-\alpha;T_0}
 \le\frac{Cr^2}{H}\qanc^{j-1}+C\sum_{i=1}^{j-1}\qanc^{j-1-i}\Bigl\{v_i\sqrt{\log(2+L/a_i)}+u_i\Bigr\}.
\]
\item For $1/2<\alpha<3/2$ and every $0<\tau<\min\{1/2,-\log\qanc/\log B\}$,
\[
 \sum_{j=l}^k\E\bigl\lVert a_{j-1}^{-1}\overline Q_{d_j^k,j}(\theta_{j-1},\theta_{j-1})\bigr\rVert_{-\alpha;T_0}\le C_\tau a_l^\tau,
\]
uniformly in $k\ge l\ge2$.
\end{enumerate}
The constants depend on the fixed construction parameters and coefficient bounds, on $L$, $H$, $T_0$, $\alpha$, and $\tau$, and on the fixed coupling $g_3$, but not on $k$, the depth, the ratio $a_i/\delta$, or the sample. The identity in the fourth part does not identify compositions of the nonlinear zero-fluctuation reconstructions with products of their linearizations; their difference is contained in $e_{j,i}$.
\end{theorem}

\begin{proof}
We fix $k$ for the whole proof, so that all coefficients, zero-fluctuation reconstructions, and conditional kernels have the depths $d_i^k=J(k-i)$ of one hierarchy. Intermediate constants in the estimates for a single kernel may depend on $J$; all estimates are uniform in the depth. The proof has five steps. Steps 1 and 2 prove the first two parts, Step 3 proves the third, Step 4 the fourth, and Step 5 the last.

\paragraph{Step 1: the conditional expectation of a quadratic term.}
Fix a level $i$ and a coarse field $e^\phi\in B_r$, and write $R=\norm\phi_\infty$, $g=g_i$, $t=t_i$, and $p=p_i$. In this step, all expectations are under the conditional fluctuation kernel at $e^\phi$. Let
\[
 \mu_\phi=\E_\phi w,\qquad\xi=w-\mu_\phi,\qquad\Sigma_\phi=\E_\phi\xi\xi^{\mathsf T},\qquad\psi_\phi(w)=\widehat\Phi_{d_i^k,i}(\phi,(g_{i,s}w_s)_s).
\]
The segment of coarse fields $u\phi$, $0\le u\le1$, stays in $B_r$. The fluctuation cutoff domain, the free logarithmic coordinates, and the conditional fluctuation potential are those preserved by \cref{ch12:thm:representative-prefixes}. The domains and bounds on the coupling-scaled fluctuation inputs are those fixed in \cref{ch14:sec:setting}; hence the whole segment satisfies the domain conditions. The bounds on means and on the dependence on the coarse field in \cref{ch10:thm:actual-means-and-scores} apply to this kernel, and give
\begin{equation}\label{ch14:eq:white-data}
 \norm{\mu_\phi}_\infty\le C_JtR,\qquad\norm{\Sigma_\phi}_{\mathrm{op}}\le2,\qquad\norm{\Sigma_\phi-\Sigma_0}_{\mathrm{row+col}}\le C_Jtp^3R,
\end{equation}
where the last norm is the larger of the maximal row and column sums, possibly with a fixed smaller exponential weight; since it is at most their sum, the cited bound holds with either convention. The product of balls does not depend on $\phi$. To make precise which normalization enters, let $V_\phi$ be the conditional fluctuation potential and $S_u=\partial_uV_{u\phi}-\E_{u\phi}\partial_uV_{u\phi}$ its centered potential derivative along the segment. In the notation of \cref{ch10:lem:fixed-domain-calculus}, $S_u=Q_u$ for this path, and $-S_u$ is the normalized log-density derivative. By differentiation, we have
\begin{equation}\label{ch14:eq:score-identity}
 \partial_u\E_{u\phi}F=\E_{u\phi}\partial_uF-\E_{u\phi}(FS_u),\qquad\partial_u\Sigma_{u\phi}(r,s)=-\E_{u\phi}(\xi_r\xi_sS_u).
\end{equation}
Indeed, we differentiate the numerator and the denominator; this is legitimate because the support is compact and fixed and the integrands are smooth. In the second identity, the terms containing the derivative of $\mu_{u\phi}$ multiply $\E\xi=0$ and cancel. Thus \eqref{ch14:eq:white-data} refers to the covariance of the interacting truncated kernel and to its centered derivative; there is no boundary term and no reference Gaussian covariance.

By the joint coefficient bounds kept by \cref{ch12:lem:deterministic-prolongation-reserve}, we may expand the reconstruction map by Taylor's formula at the conditional mean $\mu_\phi$. We define
\[
 v_\phi=\psi_\phi(\mu_\phi),\qquad E_\phi=g^{-1}D_w\psi_\phi(\mu_\phi),\qquad B_\phi=(2g^2)^{-1}D_w^2\psi_\phi(\mu_\phi);
\]
the ratios $g_{i,s}/g$ are included in these definitions. The mean lies in the convex product, and so does the segment from it to each $w$. Since $\norm\xi_\infty\le2p$, Taylor's formula with integral remainder and the bound on the third derivative in $w$ imply
\begin{equation}\label{ch14:eq:true-mean-expansion}
 \theta_i=v_\phi+gE_\phi\xi+g^2B_\phi(\xi,\xi)+\mathcal R_\phi(\xi),\qquad\norm{\mathcal R_\phi(\xi)}_\infty\le C_Jg^3p^3.
\end{equation}
Each derivative in $w$ contributes its coupling factor, and the joint absolute bound sums the remaining coefficients. These bounds are at fixed $J$ and impose no further condition on $J$.

By the bounds on the row sums and volume-normalized column sums kept by the reconstruction maps, converted between the finitely many types of fluctuation variables of a step, we have
\[
 \norm{E_\phi}_{\mathrm{op}}+\norm{B_\phi}_{\mathrm{abs}}\le C_J,
\]
where $\norm{B_\phi}_{\mathrm{abs}}$ is the larger of its absolute row and column sums. The joint bounds in the coarse field and the coupling-scaled fluctuation inputs also give
\begin{equation}\label{ch14:eq:coefficient-response}
 \norm{E_\phi-E_0}_{\mathrm{op}}+\norm{B_\phi-B_0}_{\mathrm{abs}}\le C_J(R+g\norm{\mu_\phi}_\infty)\le C_J'R.
\end{equation}
For this bound, we integrate the derivative in the coarse field along $s\phi$ with the coupling-scaled fluctuation inputs equal to zero, and then the derivative in these inputs from zero to $(g_{i,s}\mu_{\phi,s})_s$. Both segments lie in the fixed inner domain, where the normalized coefficients have bounded row and column sums. At the zero coarse field, we have $\mu_0=0$ by the global color symmetry, and the endpoint is exactly $E_0$, $B_0$. The row and column bounds imply the displayed operator bounds by Schur's test,
\[
 \sum_x\Bigl\lvert\sum_yK_{xy}u_y\Bigr\rvert^2\le\Bigl(\sup_x\sum_y\abs{K_{xy}}\Bigr)\Bigl(\sup_y\sum_x\abs{K_{xy}}\Bigr)\sum_y\abs{u_y}^2,
\]
and the conversion from volume-normalized to counting column sums costs only a factor depending on $J$.

Let $\mathcal Z_\phi=E_\phi\Sigma_\phi E_\phi^{\mathsf T}$. Adding and subtracting $E_0\Sigma_\phi E_\phi^{\mathsf T}$ and $E_0\Sigma_0E_\phi^{\mathsf T}$, and using \eqref{ch14:eq:white-data} and \eqref{ch14:eq:coefficient-response}, we obtain
\begin{equation}\label{ch14:eq:covariance-response}
 \norm{\mathcal Z_\phi-\mathcal Z_0}_{\mathrm{op}}\le C_JR(1+tp^3)\le C_J'R.
\end{equation}
This is the covariance comparison of the first part, for the coefficients of the reconstruction maps themselves.

At the zero coarse field, the product of balls and the potential are invariant under every constant adjoint color rotation. This implies $\Sigma_0=\Sigma_{\mathrm{sp}}\otimes I_3$ for a spatial matrix $\Sigma_{\mathrm{sp}}$, and the linear coefficient $E_0$ acts as a scalar on the adjoint color representation. To check this, note that rotations by $\pi$ about the three coordinate axes force each $3\times3$ color block commuting with them to be diagonal, and rotations exchanging the axes force its three diagonal entries to agree. The same invariance for the equivariant map $Q$ implies
\begin{equation}\label{ch14:eq:flat-contraction}
 Q:\mathcal Z_0=0.
\end{equation}
Indeed, each output component of the contraction is an adjoint vector fixed by every rotation, and rotations by $\pi$ about two different axes force it to vanish. This argument is componentwise in the output and keeps every spatial entry of the covariance; it needs no translation invariance of $Q$.

The same reasoning shows that $B_0:\Sigma_0=0$. Together with \eqref{ch14:eq:coefficient-response} and \eqref{ch14:eq:white-data}, this implies
\begin{equation}\label{ch14:eq:quadratic-reconstruction-mean}
 \norm{\E_\phi B_\phi(\xi,\xi)}_\infty\le C_JR(1+tp^3)\le C_J'R.
\end{equation}
Here we used that each scalar coefficient matrix of $B_\phi$ has absolute entry sum at most $C_J$, and that a contraction against a matrix of operator norm $D$ has absolute value at most $C_JD$, since every entry of that matrix is bounded by $D$. The same remark justifies contractions against \eqref{ch14:eq:covariance-response} without any dependence on the dimension.

The displacement of the reconstruction at the mean satisfies
\begin{equation}\label{ch14:eq:v-center}
 \norm{v_\phi-\psi_\phi(0)}_\infty\le C_Jg\norm{\mu_\phi}_\infty\le C_Jg^2R,
\end{equation}
where we integrated the derivative in $w$, whose absolute row sums are at most $C_Jg$, along the segment from $0$ to $\mu_\phi$, and used that $t/g$ depends only on $J$. By the bounds on the first three derivatives of the zero-fluctuation reconstruction and the majorant of \cref{ch12:lem:deterministic-prolongation-reserve}, we have
\begin{equation}\label{ch14:eq:center-tangent}
 \norm{\psi_\phi(0)-L_i\phi}_\infty\le C_Ja_{\mathrm{cap}}R^2,\qquad\norm{L_i\phi}_\infty\le a_{\mathrm{cap}}R.
\end{equation}
Only the constant in the linear bound, which is uniform in $J$, enters the choice of $a_{\mathrm{cap}}$; the quadratic bound may have a constant depending on $J$. By \eqref{ch14:eq:v-center}, \eqref{ch14:eq:center-tangent}, and $\norm{Q(u,v)}_\infty\le K_Q\norm u_\infty\norm v_\infty$, we have
\begin{equation}\label{ch14:eq:deterministic-quadratic-error}
 \norm{Q(v_\phi,v_\phi)-Q(L_i\phi,L_i\phi)}_\infty\le C_JK_Q(a_{\mathrm{cap}}^2R^3+g^2R).
\end{equation}
To see every power, let $\epsilon=v_\phi-L_i\phi$; then the left side is at most $2K_Qa_{\mathrm{cap}}R\norm\epsilon_\infty+K_Q\norm\epsilon_\infty^2$, with $\norm\epsilon_\infty\le C_J(a_{\mathrm{cap}}R^2+g^2R)$. Since $R$ and $g$ are bounded by fixed constants, the term $R^4$ is absorbed into $R^3$ and the remaining terms into $g^2R$.

We now expand $Q(\theta_i,\theta_i)$ using \eqref{ch14:eq:true-mean-expansion}. The term linear in $gE_\phi\xi$ has mean zero. The term quadratic in it has mean
\[
 \E_\phi Q(gE_\phi\xi,gE_\phi\xi)=g^2Q:\mathcal Z_\phi,
\]
which is the identity of the first part; by \eqref{ch14:eq:flat-contraction} and \eqref{ch14:eq:covariance-response}, it is at most $C_JK_Qg^2R$. The cross term of $v_\phi$ with $g^2B_\phi(\xi,\xi)$ has mean at most $C_JK_Qg^2R$, by \eqref{ch14:eq:quadratic-reconstruction-mean} and the boundedness of $v_\phi$. Every remaining product contains either one linear and one quadratic term in $\xi$, or a remainder, or two quadratic terms. Their pointwise bounds are $C_JK_Qg^3p^3$, $C_JK_Qg^3p^3$, and $C_JK_Qg^4p^4$, respectively. Products of a remainder with a term of positive degree in $\xi$ satisfy the same bound $C_JK_Qg^3p^3$, since $gp$ is bounded by a fixed constant, and so does the square of the remainder. In the term $Q(v_\phi,\mathcal R_\phi)$, only the boundedness of $v_\phi$ is needed. Further, $g^4p^4\le C_Jg^3p^3$. Combining these bounds with \eqref{ch14:eq:deterministic-quadratic-error}, we obtain the first estimate. 

\paragraph{Step 2: the centered heat norm of a quadratic term.}
We prove the second part without assuming smoothness of the output. Fix a sample, so that $\phi=\theta_{i-1}$ is fixed and $\norm\phi_\infty\le\overline b_{i-1}$. Let $Y(w)=Q(\theta_i(w),\theta_i(w))$ and $F=Y-\E_\phi Y$. On the whole product of balls, differentiating the polarized map, we find
\[
 D_wY=Q(D_w\theta_i,\theta_i)+Q(\theta_i,D_w\theta_i).
\]
The absolute row sums of the derivative of $\theta_i$ in $w$ are at most $C_Jg_i$, and its volume-normalized columns satisfy the same bound. At most a number of fluctuation variables depending only on $J$ is attached to each cell of the lattice of mesh $a_i$. Further, $\norm{\theta_i(w)}_\infty\le\overline b_i$ at every point $w$ of the product. Indeed, the recurrence of Step 1 in the proof of \cref{ch14:lem:backward-caps-and-means} gives $\norm{\theta_i(w)}_\infty\le a_{\mathrm{cap}}\overline b_{i-1}+C_{\mathrm{rec},J}t_ip_i$, and this is at most $\overline b_i$ by \eqref{ch14:eq:cap-constants}, since $C_J\ge C_{\mathrm{cap},J}$. Let $\mathcal A=C_JK_Qg_i\overline b_i$. Then the derivative matrix $G=D_wY$ satisfies
\begin{equation}\label{ch14:eq:quadratic-jacobian}
 \sup_y\sum_r\abs{G_{y,r}}\le\mathcal A,\qquad\sup_r\frac{\delta^3}{a_i^3}\sum_y\abs{G_{y,r}}\le\mathcal A,
\end{equation}
where $r$ runs over the scalar fluctuation variables. For the column bound, we fix the differentiated fluctuation variable, use the column of $Q$ for its differentiated argument, bound its other argument by $\overline b_i$, and then use the volume-normalized column of $D_w\theta_i$; the volume factors multiply to $\delta^3/a_i^3$. Since both differentiated arguments satisfy the same bound, their sum is included in $C_J$. This holds for all directions and colors.

By Schur's test, the counting $\ell^2$ operator norm of $G$ is at most $\mathcal A(a_i/\delta)^{3/2}$. For $0<s\le T_0$, let $\omega_y(s,x)=\int_{Q_y}p_s(x-z)\,dz$ for the output cells $Q_y$ of side $\delta$. For each scalar component of $P_sY$ at a point $x$, the Lipschitz constant in $w$ is at most
\begin{equation}\label{ch14:eq:quadratic-heat-row}
 C\mathcal A\min\{1,a_i^{3/2}s^{-3/4}\}.
\end{equation}
The first bound follows from the absolute row sums and $\sum_y\omega_y=1$. For the second, we compose the operator bound with $\norm{\omega(s,x)}_2\le C\delta^{3/2}s^{-3/4}$, from \cref{ch13:lem:heat-and-chaining}; the factor $(a_i/\delta)^{3/2}$ cancels the factor $\delta^{3/2}$ of the output cells. By the derivative and time-increment bounds of the same lemma, one spatial derivative of the cell integrals costs $s^{-1/2}$ and one time derivative costs $s^{-1}$ when $a_i^2\le s\le T_0$, in which case $s\ge\delta^2$. These bounds also hold for differences of two evaluations, by integration along a shortest spatial path or a time segment.

We now take the supremum over times, beginning with times below $a_i^2$. Let $M_i=\E_\phi\theta_i$. By the amplitude bound, we have $\norm{M_i}_\infty\le\overline b_i$ as well as $\norm{\theta_i}_\infty\le\overline b_i$. The output row bound in \eqref{ch14:eq:quadratic-norm} then gives
\[
 \norm{Q(\theta_i,\theta_i)-Q(M_i,M_i)}_\infty\le2K_Q\overline b_i\norm{\theta_i-M_i}_\infty,
\]
and subtracting the conditional expectation, we obtain
\begin{equation}\label{ch14:eq:native-quadratic-deviation}
 \norm F_\infty\le2K_Q\overline b_i\bigl(\norm{\theta_i-M_i}_\infty+\E_\phi\norm{\theta_i-M_i}_\infty\bigr).
\end{equation}
This uses the output row; a bound on the input columns alone would not imply it. There are $N_i=9(L/a_i)^3$ scalar coordinates of $\theta_i$ on its own lattice. By the bounds on the derivative rows of the reconstruction maps in \cref{ch12:lem:deterministic-prolongation-reserve}, each of them is a Lipschitz function of $w$ with Euclidean constant at most $C_Jg_i$, since the $\ell^2$ norm of a row of the derivative is at most its absolute sum. The conditional kernel is the one on a convex product described in \cref{ch14:sec:setting} (\cref{ch12:lem:backward-law,ch12:thm:representative-prefixes}), with Hessian at least $I/2$. By \cref{ch02:thm:reflecting-concentration} with $m=1/2$ and a union bound, we have
\begin{equation}\label{ch14:eq:native-coordinate-maximum}
 \Prob_\phi\{\norm{\theta_i-M_i}_\infty>x\}\le\min\bigl\{1,2N_i\exp[-x^2/(C_Jg_i^2)]\bigr\}.
\end{equation}
To integrate this bound, we take $x_0=g_i\sqrt{C_J\log(2N_i)}$, and for this integral we let $A^2=C_Jg_i^2$. We use $\int_{x_0}^\infty e^{-x^2/A^2}\,dx\le A^2e^{-x_0^2/A^2}/(2x_0)$, which follows by inserting the factor $x/x_0\ge1$ in the integrand. After enlarging the constant, we obtain
\[
 \E_\phi\norm{\theta_i-M_i}_\infty\le C_Jg_i\sqrt{\log(2+L/a_i)},\qquad\E_\phi\norm F_\infty\le C_JK_Qg_i\overline b_i\sqrt{\log(2+L/a_i)},
\]
where we used $\log(2N_i)\le C\log(2+L/a_i)$, which holds because $L/a_i\ge4$. Since the heat semigroup contracts the supremum norm, the expected supremum over $0<s\le\min\{a_i^2,T_0\}$ of $s^{\alpha/2}\norm{P_sF}_\infty$ is at most $C\mathcal Aa_i^\alpha\sqrt{\log(2+L/a_i)}$. Only the coordinates on the lattice of mesh $a_i$ have been counted; there is no union bound over the arbitrarily fine output lattice. The mean $M_i$ is the conditional mean of the reconstruction, which need not vanish.

For the remaining times, let $S_h=2^ha_i^2$, and consider each nonempty interval $[S_h,\min\{2S_h,T_0\}]$, $h\ge0$. The centered process $s^{\alpha/2}(P_sF)(x)$, with its nine scalar components, has one-point sub-Gaussian parameter at most
\[
 \sigma_h=C_{J,\alpha}\mathcal Aa_i^{3/2}S_h^{\alpha/2-3/4},
\]
by \eqref{ch14:eq:quadratic-heat-row} and \eqref{ch13:eq:shell-product-inequalities}. Its increments have parameter at most $C\sigma_h$ times
\[
 \min\{1,\dist(x,x')/\sqrt{S_h}\}+\min\{1,\abs{s-s'}/S_h\};
\]
this includes the derivative of $s^{\alpha/2}$, whose relative size on the interval is at most $C_\alpha/S_h$. Between distinct components, we take the distance to be $C\sigma_h$; adding the Lipschitz constants of the two components justifies this choice. Using spatial nets of spacing $c\eta\sqrt{S_h}$ and temporal nets of spacing $c\eta S_h$, together with the nine components, we obtain covering numbers at most
\[
 C(1+L/\sqrt{S_h})^3\eta^{-4}\qquad(0<\eta<1).
\]
Since the process is continuous for $s>0$ and smooth in $x$ on these intervals, a countable dense set of parameters determines its supremum. We apply \cref{ch13:lem:coefficient-metric-chaining} to the increment function divided by $\sigma_h$, with $s=C\sigma_h$, $b=0$, $D=C$, $A=C(1+L/\sqrt{S_h})^3$, and $q_{\mathrm{ent}}=4$; the point $z_0$ of that lemma has the same one-point bound. At $p=1$, the lemma bounds the expected supremum over the interval by
\[
 C\sigma_h\sqrt{\log(2+L/a_i)},
\]
because $\sqrt{S_h}\ge a_i$, so that $A\le C(1+L/a_i)^3$; no factor depending on the output mesh enters. Since $\sigma_h=C_{J,\alpha}\mathcal Aa_i^\alpha2^{-h(3/4-\alpha/2)}$, the sum over the intervals converges with ratio $2^{\alpha/2-3/4}<1$ when $\alpha<3/2$. Adding the bound for small times, we obtain
\[
 \E_\phi\norm F_{-\alpha;T_0}\le C_{J,\alpha}K_Qg_i\overline b_ia_i^\alpha\sqrt{\log(2+L/a_i)}.
\]
Conditionally on $\mathcal F_{i-1}$, all constants here are deterministic and uniform over the coarse fields that occur. This is the second part.

\paragraph{Step 3: the pulled-back coefficients.}
The linearized reconstruction $L_i$ has output mesh $a_i$ and input mesh $a_{i-1}=Ba_i$, and it satisfies \eqref{ch14:eq:tangent-rows-columns}. These are the weighted bounds on the first derivative of the zero-fluctuation reconstruction, with the constant $C_{\mathrm{CP}}$, uniform in $J$, enlarged as in \cref{ch14:sec:setting}. By the bounds on the second derivative of the zero-fluctuation reconstruction, the coefficient $Q_{j,j-1}=\overline Q_{d_j^k,j}$ has norm at most a fixed constant $C$. By induction down to the input level $i$, using \cref{ch14:lem:quadratic-pullback} at each step, we obtain
\[
 K_{j,i}\le Ca_{\mathrm{cap}}^{2(j-1-i)}.
\]
The global color equivariance of the zero-fluctuation reconstruction implies the equivariance of its quadratic part, and each $L_i$ commutes with constant color rotations; hence every $Q_{j,i}$ is equivariant. With $m=j-1-i$, the identity $a_i=B^ma_{j-1}$ gives
\begin{equation}\label{ch14:eq:physical-contraction}
 \frac{K_{j,i}}{a_{j-1}}\le\frac{CB^ma_{\mathrm{cap}}^{2m}}{a_i}=\frac{C\qanc^m}{a_i}.
\end{equation}

\paragraph{Step 4: the telescoping identity and the convolution.}
Let $Y_i=Q_{j,i}(\theta_i,\theta_i)$ for $0\le i\le j-1$. By the definitions in the fourth part,
\[
 Y_i-Y_{i-1}=\bigl\{Y_i-\E[Y_i\mid\mathcal F_{i-1}]\bigr\}+\bigl\{\E[Y_i\mid\mathcal F_{i-1}]-Y_{i-1}\bigr\}=Z_{j,i}+e_{j,i}.
\]
Summing this finite identity over $1\le i\le j-1$, we obtain the telescoping identity on every sample. Since every conditional expectation is taken under the kernel of the same hierarchy that generated $\theta_i$, Steps 1 and 2 apply with $\phi=\theta_{i-1}$ and the amplitude bound $\overline b_{i-1}$ on $\norm{\theta_{i-1}}_\infty$.

The term from the coarsest level is bounded deterministically by
\[
 \Bigl\lVert\frac{Q_{j,0}(\theta_0,\theta_0)}{a_{j-1}}\Bigr\rVert_{-\alpha;T_0}\le\frac{Cr^2}{H}\qanc^{j-1},
\]
by \eqref{ch14:eq:bounded-heat}, \eqref{ch14:eq:physical-contraction} with $i=0$, $a_0=H$, and $\norm{\theta_0}_\infty\le r$. For the centered terms, the second part and \eqref{ch14:eq:physical-contraction} give
\[
 \frac{\E\norm{Z_{j,i}}_{-\alpha;T_0}}{a_{j-1}}\le C\qanc^{j-1-i}g_i\overline b_ia_i^{\alpha-1}\sqrt{\log(2+L/a_i)}.
\]
For the conditional expectation errors, we use the first part, \eqref{ch14:eq:bounded-heat}, and the three identities
\[
 \frac{a_{\mathrm{cap}}^2}{a_i}=\frac{\qanc}{a_{i-1}},\qquad\frac{g_i^2}{a_i}=g_3^2,\qquad\frac{g_i^3}{a_i}=g_3^3\sqrt{a_i},
\]
and obtain
\[
 \frac{\E\norm{e_{j,i}}_{-\alpha;T_0}}{a_{j-1}}\le C\qanc^{j-1-i}u_i.
\]
By the triangle inequality, we obtain the convolution estimate of the fourth part, with $v_i$ as defined there. No composition of the nonlinear zero-fluctuation reconstructions has been identified with a product of linearizations; the difference remains in $e_{j,i}$.

\paragraph{Step 5: summation.}
Let $1/2<\alpha<3/2$. By the reconstruction amplitude bound of \cref{ch14:lem:backward-caps-and-means} and \eqref{ch14:eq:p-affine}, we have, with constants depending on the fixed $H$, $g_3$, and $r$,
\begin{equation}\label{ch14:eq:cap-decay}
 \overline b_i\le C(1+i)B^{-i/2}\qquad(i\ge0),
\end{equation}
since $a_{\mathrm{cap}}^i=B^{-\kapterm i}$ with $\kapterm>1/2$; the case $i=0$ is included by enlarging $C$. More precisely, inserting the two terms of $\overline b_i$, we obtain
\[
 v_i\le C\bigl\{B^{-(\kapterm+\alpha-1/2)i}+(1+i)B^{-\alpha i}\bigr\},
\]
and both exponents exceed $1/2$. In $u_i$, the part of $\overline b_{i-1}^3/a_{i-1}$ coming from the coarsest field decays with the exponent $3\kapterm-1>1/2$, and the part coming from the fluctuations is at most $C(1+i)^3B^{-i/2}$; the inequality $(x+y)^3\le4(x^3+y^3)$ handles the mixed terms. The other two terms of $u_i$ are at most $C(1+i)^3B^{-i/2}$ directly. It follows that
\begin{equation}\label{ch14:eq:vu-decay}
 v_i+u_i\le C(1+i)^3B^{-i/2}.
\end{equation}
For $i<j$, we use the cruder bound $\sqrt{\log(2+L/a_i)}\le\sqrt{\log(2+L/a_j)}\le C\sqrt{1+j}$.

Let $0<\tau<\min\{1/2,-\log\qanc/\log B\}$, and set $r_1=\qanc B^\tau<1$ and $r_2=B^{\tau-1/2}<1$. Let $\mathcal X_j$ be the deterministic upper bound obtained from the convolution estimate by inserting \eqref{ch14:eq:vu-decay} and $\sqrt{\log(2+L/a_i)}\le C\sqrt{1+j}$; it is defined for every $j\ge2$ and does not depend on $k$. Writing $j=i+m+1$, we have
\[
 \sum_{j\ge2}B^{\tau j}\mathcal X_j\le C\sum_{j\ge2}B^{\tau j}\qanc^{j-1}+CB^\tau\sum_{i\ge1}\sum_{m\ge0}(2+i+m)^{1/2}(1+i)^3r_1^mr_2^i<\infty.
\]
All terms are nonnegative, and $(2+i+m)^{1/2}\le(2+i)(1+m)$ reduces the double series to the product of $\sum_i(2+i)^4r_2^i$ and $\sum_m(1+m)r_1^m$, which converge by comparison with geometric series. The bound depends only on the fixed data and on the gaps in the strict inequalities for $\tau$. The finite sequence for a given $k$ is bounded by this deterministic majorant; no coefficients from different values of $k$ are identified. We conclude that
\[
 \sum_{j=l}^k\E\bigl\lVert a_{j-1}^{-1}\overline Q_{d_j^k,j}(\theta_{j-1},\theta_{j-1})\bigr\rVert_{-\alpha;T_0}\le B^{-\tau l}\sum_{j\ge2}B^{\tau j}\mathcal X_j\le C_\tau H^{-\tau}a_l^\tau,
\]
and the factor $H^{-\tau}$ is absorbed into $C_\tau$. This is the last part. The proof used only $\qanc<1$ and the stated choice of $\tau$; no condition of the form $\qanc B^\alpha<1$ is needed.
\end{proof}
